\subsection*{Build and Compile}

Croqtile files use the \texttt{.co} extension. The compiler works like \texttt{gcc} or \texttt{clang}:

\begin{lstlisting}[language={},caption={Build and Compile}]
croqtile ele_add.co                          # produces a.out
croqtile ele_add.co -o ele_add               # specify output name
croqtile -es -t target ele_add.co -o out.cu    # emit target source only
\end{lstlisting}

The \texttt{-t} flag selects the target platform.

To see all available options, run \texttt{croqtile -\allowbreak{}-\allowbreak{}help}:

\begin{lstlisting}[language={},caption={Build and Compile}]
croqtile --help
\end{lstlisting}

This prints the full set of compiler flags --- output naming, target selection, verbosity, and diagnostic options. It is worth scanning once so you know what is available when you need it.

\begin{figure}[htbp]\centering
\includegraphics[width=0.86\linewidth,height=0.42\textheight,keepaspectratio]{f03-compile_and_run_light.png}
\end{figure}
\begin{center}\small\emph{Compiling and running a Croqtile program, plus the \texttt{--help} output.}\end{center}
